\section{Introduction}
% ================================================================================

AI agents increasingly operate as long-lived processes that outlive the
interactive client that initiated them. A mobile user asks a research agent
to summarise a dataset; the agent enqueues several tool calls and then
pauses for a human confirmation. The user's phone moves between cell towers,
the transport disconnects, and the user does not return for forty minutes.
During that interval, what must the bridge component that mediates between
the client and the agent runtime actually \emph{do}? Existing agent
protocols answer the questions of identity (AGTP), audit (GAR) and
resumption (PACE) but leave the bridge's own question orphaned. AGTP \cite{I-D.hood-agtp} and its session companion
\cite{I-D.hood-agtp-session} define transport-level identity headers
(\texttt{Agent-ID}, \texttt{Owner-ID}, \texttt{Authority-Scope},
\texttt{Session-ID}) but do not specify the execution semantics of the
bridge that holds them. GAR \cite{I-D.sato-soos-gar} defines a tamper-evident
audit record (SAR) and the Authority Lifecycle Event (ALE), but presupposes
the existence of a Governing Enforcement Component (GEC) and does not say
who emits audit events when no GEC is present. PACE
\cite{I-D.nandakumar-agentproto-moq-pace} defines a re-attach mechanism for
sessions in the SUSPENDED state but presumes a relay or Media over QUIC
Transport (MOQT) substrate and does not define what happens in the window
\emph{during} which the client is detached and the agent continues to work.
The orphaned question---what the bridge must do during this window---is
unanswered, in the versions read, by AGTP, GAR or PACE.

This document answers it. We define the \emph{Bridge Contract}, a
structured specification of the tuple $BC = (S, E, R, C)$ where $S$ is a
Session Identity Descriptor, $E$ is the set of Governance Events the
bridge is obligated to emit, $R$ is the Resume Record the
bridge keeps for itself on detach, and $C$ is the set of Conformance
Criteria against which any implementation can be independently
verified. The contract is positioned as \emph{execution semantics for
AGTP headers}, \emph{a GAR-shaped event source for deployments
without a Governing Enforcement Component (GEC)}, and \emph{the
source of the resume point that a PACE re-attach
needs}. It does not compete with any of these protocols.

The research questions RQ1--RQ4 answered in this document are:

\begin{description}[leftmargin=2em]
\item[RQ1] Given that AGTP defines transport headers such as
\texttt{Session-ID}, \texttt{Authority-Scope}, and \texttt{Task-ID}, are
there uncovered governance obligations at the bridge layer? If so, which?
\item[RQ2] GAR already defines the Session Audit Record (SAR) and the
Authority Lifecycle Event (ALE). Is the Bridge Contract a competitor, a
complement, or an ``execution-side precondition'' for GAR?
\item[RQ3] PACE already defines a session resumption mechanism (re-attach
to SUSPENDED sessions). Does the Bridge Contract's resume record and re-attach protocol
overlap with PACE, complement it, or duplicate it? If it complements, where
does the boundary lie?
\item[RQ4] What is the Bridge Contract's uniqueness with respect to
existing session-level implementations such as Herdr? What exactly is
contributed, and what is the evidence for it?
\end{description}

\paragraph{Terminology.} In this document, ``Bridge Contract'' (BC) refers to
the specification defined hereby; ``bridge'' refers to the component that
implements the contract. ``Resume Record'' (Definition~\ref{def:rt})
is the bridge's own record of where a detachment can be resumed; it is
not handed to the client. The client re-attaches by challenge-response
(Section~\ref{sec:spec}); a MACed ``resume token'' held by the client
was considered and rejected, for the reason given in
Section~\ref{sec:spec}.
``GEC-less deployment'' refers to a deployment without
a Governing Enforcement Component, per Definition~\ref{def:gecless}.

\paragraph{Novelty statement.} The contribution of this document is one
artifact with a specification, evidence and a gap analysis: a \emph{specification} (a tuple, a normative
state machine and conformance criteria) for the obligations that no
protocol among AGTP, GAR, and PACE assigns to the bridge, \emph{executable
evidence} (reference implementations, tests, mutants and vectors), and a
\emph{formal model} checked within stated bounds. The \emph{phenomenon} of orphaned obligations in
protocol composition is reported by AgentThread
\cite{zheng2026formal} for other protocols, which motivates this work. Hash-chained
audit events, capability revocation during detach, and event-before-mutation
audit patterns are standard techniques in the protocol composition
literature; the Bridge Contract does not introduce new mechanisms in any
of these areas. We use the term ``Bridge Contract'' in the engineering
sense (a structurally specified set of obligations), not in the sense of
session types or behavioural contracts.

\paragraph{Contributions.} The contributions of this document are:
\begin{enumerate}[leftmargin=2em]
\item A structured specification of the Bridge Contract as a tuple
$BC = (S, E, R, C)$ with formal definitions of $S$, $E$, $R$, $C$
(Section~\ref{sec:spec}), where $R$ is a resume record the bridge keeps
and re-attach is a challenge-response over a single-use nonce and the
last event the client saw. According to the author's report on the Termenbox integration, a
client-held MACed token cannot work (Section~\ref{sec:impl}).
\item A normative product state machine
$\{$attachment$\} \times \{$human$\} \times \{$lifecycle$\}$ with
allowed transitions, preconditions and events, restart recovery, several
clients per principal (attachment derived from the client set, I15), and
conformance invariants I1--I15, among them proof of possession over a
single-use nonce (I11), durable append (I12), queue integrity (I13) and
answer binding (I14) (Section~\ref{sec:state-machine}). The state
names ATTACHED, DETACHED, PENDING\_HUMAN, RESUMING and TERMINATED remain
only as a derived view. According to the author's report, the integration motivated I15 and the client set
(a single attached client is not enough), I14, and I9, which bounds a
nonce rather than a token.
\item Explicit axioms (A1--A7), lemmas
(Lemmas~1--4), and liveness properties (L1--L2) supporting the safety
argument (Section~\ref{sec:state-machine}), with timers derived from
one value, $T_{\text{detach}}$, with $T_{\text{call}} \le T_{\text{human}}$ as a
bound and $T_{\text{human}} \le T_{\text{detach}}$ as the default,
and two profiles; as reported by the author, the integration motivated the derivation.
\item A statement of what the bridge can see (Section~\ref{sec:coverage}):
I2 and I13 cover calls the runtime submits to the bridge, and I3 splits
into scope widening the bridge prevents and widening the runtime
announces, which it reports.
\item A comparison of each Bridge Contract obligation with the cited drafts
and one product, with its method and a generated matrix of cells
(Section~\ref{sec:mapping}).
\item A reference implementation note describing how a bridge can satisfy
the contract without requiring changes to AGTP, GAR, or PACE
(Section~\ref{sec:impl}).
\item Evidence that the specification is free of invariant violations on bounded models and is testable
(Section~\ref{sec:conformance}, Appendix~\ref{app:formal}): a TLC-checked
TLA+ model of the product machine with a mutant check, a Python
reference bridge with a pytest suite (including
property-based tests) and a mutation check with a killed mutant for
every invariant I1--I15, language-neutral JSON
test vectors V01--V134, and a second implementation in Rust, behind storage and
verifier traits, that runs the same vectors except V52--V54, V88 and V89, which expect
checkpoint events the crate does not emit. Both implementations have the same author; the evidence shows
consistency, not independent confirmation.
\end{enumerate}

\paragraph{Roadmap.} The remainder of this document is structured as
follows. Section~\ref{sec:spec} defines the Bridge Contract tuple
$BC = (S, E, R, C)$. Section~\ref{sec:state-machine} specifies the
state machine, invariants, axioms, lemmas, liveness, timeouts, and
the safety argument. Section~\ref{sec:mapping} compares each obligation with
the cited drafts and presents the threat model.
Section~\ref{sec:impl} sketches a reference implementation.
Section~\ref{sec:conformance} presents the conformance test suite.
Section~\ref{sec:limitations} states limitations and out-of-scope
items. Appendix~\ref{app:formal} provides the TLA+ model of the product
machine and its TLC results.

% ================================================================================
